\documentclass[10pt,a4paper]{amsart}
\usepackage{amsmath,amssymb,amsthm,graphicx}
% Retained minimal preamble for the Stacks Project chapter more-morphisms.tex (More on Morphisms).
% Verbatim excerpt of the chapter preamble: Xy-pic setup (native preamble lines 14-18) and the native
% theorem declarations with the mathematical macros (native preamble lines 154-217). Copied mechanically;
% nothing added, removed or redefined. The wrapper class and the amsmath/amssymb/amsthm/graphicx base
% packages are supplied by the assembler, so the class/externaldocument part of the native preamble is omitted.
% For commutative diagrams we use Xy-pic
\usepackage[all]{xy}

% We use 2cell for 2-commutative diagrams.
\xyoption{2cell}
% Theorem environments.
%
\theoremstyle{plain}
\newtheorem{theorem}[subsection]{Theorem}
\newtheorem{proposition}[subsection]{Proposition}
\newtheorem{lemma}[subsection]{Lemma}

\theoremstyle{definition}
\newtheorem{definition}[subsection]{Definition}
\newtheorem{example}[subsection]{Example}
\newtheorem{exercise}[subsection]{Exercise}
\newtheorem{situation}[subsection]{Situation}

\theoremstyle{remark}
\newtheorem{remark}[subsection]{Remark}
\newtheorem{remarks}[subsection]{Remarks}

\numberwithin{equation}{subsection}

% Macros
%
\def\lim{\mathop{\mathrm{lim}}\nolimits}
\def\colim{\mathop{\mathrm{colim}}\nolimits}
\def\Spec{\mathop{\mathrm{Spec}}}
\def\Hom{\mathop{\mathrm{Hom}}\nolimits}
\def\Ext{\mathop{\mathrm{Ext}}\nolimits}
\def\SheafHom{\mathop{\mathcal{H}\!\mathit{om}}\nolimits}
\def\SheafExt{\mathop{\mathcal{E}\!\mathit{xt}}\nolimits}
\def\Sch{\mathit{Sch}}
\def\Mor{\mathop{\mathrm{Mor}}\nolimits}
\def\Ob{\mathop{\mathrm{Ob}}\nolimits}
\def\Sh{\mathop{\mathit{Sh}}\nolimits}
\def\NL{\mathop{N\!L}\nolimits}
\def\CH{\mathop{\mathrm{CH}}\nolimits}
\def\proetale{{pro\text{-}\acute{e}tale}}
\def\etale{{\acute{e}tale}}
\def\QCoh{\mathit{QCoh}}
\def\Ker{\mathop{\mathrm{Ker}}}
\def\Im{\mathop{\mathrm{Im}}}
\def\Coker{\mathop{\mathrm{Coker}}}
\def\Coim{\mathop{\mathrm{Coim}}}

% Boxtimes
%
\DeclareMathSymbol{\boxtimes}{\mathbin}{AMSa}{"02}

%
% Macros for moduli stacks/spaces
%
\def\QCohstack{\mathcal{QC}\!\mathit{oh}}
\def\Cohstack{\mathcal{C}\!\mathit{oh}}
\def\Spacesstack{\mathcal{S}\!\mathit{paces}}
\def\Quotfunctor{\mathrm{Quot}}
\def\Hilbfunctor{\mathrm{Hilb}}
\def\Curvesstack{\mathcal{C}\!\mathit{urves}}
\def\Polarizedstack{\mathcal{P}\!\mathit{olarized}}
\def\Complexesstack{\mathcal{C}\!\mathit{omplexes}}
% \Pic is the operator that assigns to X its picard group, usage \Pic(X)
% \Picardstack_{X/B} denotes the Picard stack of X over B
% \Picardfunctor_{X/B} denotes the Picard functor of X over B
\def\Pic{\mathop{\mathrm{Pic}}\nolimits}
\def\Picardstack{\mathcal{P}\!\mathit{ic}}
\def\Picardfunctor{\mathrm{Pic}}
\def\Deformationcategory{\mathcal{D}\!\mathit{ef}}
\newenvironment{reference}{}{}
\newenvironment{slogan}{}{}
\newenvironment{history}{}{}
\usepackage{cleveref}
\crefname{originallemma}{original lemma}{original lemmas}
\Crefname{originallemma}{Original lemma}{Original lemmas}
\makeatletter
\newlabel{morphisms-lemma-H-quasi-projective-open-H-projective}{{01WA}{source}{}{sourcepointer-1}{}}
\newlabel{morphisms-lemma-H-quasi-projective-open-H-projective@cref}{{[originallemma][0][]01WA}{[1][0][]source}{}{}{}}
\makeatother
\crefname{originallemma}{original lemma}{original lemmas}
\Crefname{originallemma}{Original lemma}{Original lemmas}
\makeatletter
\newlabel{morphisms-lemma-H-quasi-projective-quasi-projective}{{01VY}{source}{}{sourcepointer-2}{}}
\newlabel{morphisms-lemma-H-quasi-projective-quasi-projective@cref}{{[originallemma][0][]01VY}{[1][0][]source}{}{}{}}
\makeatother
\crefname{originallemma}{original lemma}{original lemmas}
\Crefname{originallemma}{Original lemma}{Original lemmas}
\makeatletter
\newlabel{morphisms-lemma-ample-very-ample}{{01VN}{source}{}{sourcepointer-3}{}}
\newlabel{morphisms-lemma-ample-very-ample@cref}{{[originallemma][0][]01VN}{[1][0][]source}{}{}{}}
\makeatother
\crefname{originallemma}{original lemma}{original lemmas}
\Crefname{originallemma}{Original lemma}{Original lemmas}
\makeatletter
\newlabel{morphisms-lemma-finite-type-ample-very-ample}{{01VU}{source}{}{sourcepointer-4}{}}
\newlabel{morphisms-lemma-finite-type-ample-very-ample@cref}{{[originallemma][0][]01VU}{[1][0][]source}{}{}{}}
\makeatother
\crefname{originallemma}{original lemma}{original lemmas}
\Crefname{originallemma}{Original lemma}{Original lemmas}
\makeatletter
\newlabel{morphisms-lemma-projective-over-quasi-projective-is-H-projective}{{087S}{source}{}{sourcepointer-5}{}}
\newlabel{morphisms-lemma-projective-over-quasi-projective-is-H-projective@cref}{{[originallemma][0][]087S}{[1][0][]source}{}{}{}}
\makeatother
\crefname{originallemma}{original lemma}{original lemmas}
\Crefname{originallemma}{Original lemma}{Original lemmas}
\makeatletter
\newlabel{morphisms-lemma-pullback-ample-tensor-relatively-ample}{{0892}{source}{}{sourcepointer-6}{}}
\newlabel{morphisms-lemma-pullback-ample-tensor-relatively-ample@cref}{{[originallemma][0][]0892}{[1][0][]source}{}{}{}}
\makeatother
\crefname{originallemma}{original lemma}{original lemmas}
\Crefname{originallemma}{Original lemma}{Original lemmas}
\makeatletter
\newlabel{morphisms-lemma-quasi-projective-finite-type-over-S}{{01VS}{source}{}{sourcepointer-7}{}}
\newlabel{morphisms-lemma-quasi-projective-finite-type-over-S@cref}{{[originallemma][0][]01VS}{[1][0][]source}{}{}{}}
\makeatother
\title{Five equivalent descriptions of quasi-projective morphisms over a base with an ample invertible sheaf}
\author{The Stacks Project Authors}
\date{stacks-project:more-morphisms.tex:local-snapshot-20261011; GNU FDL 1.2 or later}
\begin{document}
\maketitle
\noindent\textit{Editorial note.} Statement span is lines 13839-13854 of sources/stacks-more-morphisms/pinned-original-source.tex; proof span is lines 13856-13889, the complete author proof (32 body lines).\par\medskip
\noindent\textit{Editorial note.} Navigation stubs: 7 cited-but-not-included labels are supplied as non-load-bearing pointers with their pinned tags where the snapshot tag index has them.\par\medskip
\noindent\textit{Editorial note.} extra\_packages is empty on purpose: the retained preamble already loads xy with options and a second declaration clashes in the assembler wrapper.\par\medskip
\noindent\textit{Editorial note.} Text is retained verbatim, including the author's own typos and placeholders.\par\medskip
\noindent\textit{Editorial note.} Hub lemma: several later chapter proofs cite this equivalence; collected for the explicit (3) implies (4)/(5) construction.\par\medskip
\noindent\textit{Allowed standard prerequisites (editorial).}\begin{itemize}
\item Explicit prerequisite boundary. The Morphisms chapter dictionary between quasi-projective, H-quasi-projective, open-H-projective and very-ample descriptions (lemmas H-quasi-projective-quasi-projective, projective-over-quasi-projective-is-H-projective, H-quasi-projective-open-H-projective, quasi-projective-finite-type-over-S, ample-very-ample, pullback-ample-tensor-relatively-ample, finite-type-ample-very-ample).
\end{itemize}\medskip
\ifdefined\hypertarget\hypertarget{sourcepointer-1}{}\fi
\noindent\textit{Source reference.} Original lemma 01WA: sources/stacks-more-morphisms/pinned-original-source.tex, source lines 13889--13889. Cited by the retained statement or proof but not part of the collected spans; supplied only as a navigation pointer. The argument does not depend on its content.\par\medskip
\ifdefined\hypertarget\hypertarget{sourcepointer-2}{}\fi
\noindent\textit{Source reference.} Original lemma 01VY: sources/stacks-more-morphisms/pinned-original-source.tex, source lines 13889--13889. Cited by the retained statement or proof but not part of the collected spans; supplied only as a navigation pointer. The argument does not depend on its content.\par\medskip
\ifdefined\hypertarget\hypertarget{sourcepointer-3}{}\fi
\noindent\textit{Source reference.} Original lemma 01VN: sources/stacks-more-morphisms/pinned-original-source.tex, source lines 13889--13889. Cited by the retained statement or proof but not part of the collected spans; supplied only as a navigation pointer. The argument does not depend on its content.\par\medskip
\ifdefined\hypertarget\hypertarget{sourcepointer-4}{}\fi
\noindent\textit{Source reference.} Original lemma 01VU: sources/stacks-more-morphisms/pinned-original-source.tex, source lines 13889--13889. Cited by the retained statement or proof but not part of the collected spans; supplied only as a navigation pointer. The argument does not depend on its content.\par\medskip
\ifdefined\hypertarget\hypertarget{sourcepointer-5}{}\fi
\noindent\textit{Source reference.} Original lemma 087S: sources/stacks-more-morphisms/pinned-original-source.tex, source lines 13889--13889. Cited by the retained statement or proof but not part of the collected spans; supplied only as a navigation pointer. The argument does not depend on its content.\par\medskip
\ifdefined\hypertarget\hypertarget{sourcepointer-6}{}\fi
\noindent\textit{Source reference.} Original lemma 0892: sources/stacks-more-morphisms/pinned-original-source.tex, source lines 13889--13889. Cited by the retained statement or proof but not part of the collected spans; supplied only as a navigation pointer. The argument does not depend on its content.\par\medskip
\ifdefined\hypertarget\hypertarget{sourcepointer-7}{}\fi
\noindent\textit{Source reference.} Original lemma 01VS: sources/stacks-more-morphisms/pinned-original-source.tex, source lines 13889--13889. Cited by the retained statement or proof but not part of the collected spans; supplied only as a navigation pointer. The argument does not depend on its content.\par\medskip
\begin{lemma}
\label{lemma-quasi-projective}
Let $S$ be a scheme which has an ample invertible sheaf.
Let $f : X \to S$ be a morphism of schemes. The following are
equivalent
\begin{enumerate}
\item $X \to S$ is quasi-projective,
\item $X \to S$ is H-quasi-projective,
\item there exists a quasi-compact open immersion $X \to X'$ of schemes
over $S$ with $X' \to S$ projective,
\item $X \to S$ is of finite type and $X$ has an ample invertible
sheaf, and
\item $X \to S$ is of finite type and there exists an
$f$-very ample invertible sheaf.
\end{enumerate}
\end{lemma}

\begin{proof}
The implication (2) $\Rightarrow$ (1) is
Morphisms, Lemma \ref{morphisms-lemma-H-quasi-projective-quasi-projective}.
The implication (1) $\Rightarrow$ (2) is
Morphisms, Lemma
\ref{morphisms-lemma-projective-over-quasi-projective-is-H-projective}.
The implication (2) $\Rightarrow$ (3) is
Morphisms, Lemma \ref{morphisms-lemma-H-quasi-projective-open-H-projective}

\medskip\noindent
Assume $X \subset X'$ is as in (3). In particular $X \to S$ is
of finite type. By
Morphisms, Lemma \ref{morphisms-lemma-H-quasi-projective-open-H-projective}
the morphism $X \to S$ is H-projective.
Thus there exists a quasi-compact immersion $i : X \to \mathbf{P}^n_S$.
Hence $\mathcal{L} = i^*\mathcal{O}_{\mathbf{P}^n_S}(1)$
is $f$-very ample. As $X \to S$ is quasi-compact we conclude from
Morphisms, Lemma \ref{morphisms-lemma-ample-very-ample}
that $\mathcal{L}$ is $f$-ample. Thus $X \to S$ is quasi-projective
by definition.

\medskip\noindent
The implication (4) $\Rightarrow$ (2) is
Morphisms, Lemma \ref{morphisms-lemma-quasi-projective-finite-type-over-S}.

\medskip\noindent
The implication (5) $\Rightarrow$ (4) follows from
Morphisms, Lemma \ref{morphisms-lemma-pullback-ample-tensor-relatively-ample}.

\medskip\noindent
The equivalence of (1) and (5) follows from
Morphisms, Lemmas \ref{morphisms-lemma-ample-very-ample} and
\ref{morphisms-lemma-finite-type-ample-very-ample}.
\end{proof}

\end{document}
